\section{Addition Accounting}
\label{sec:addition-count}

The preceding theorem optimizes the expensive gates, but the displayed
circuits also give a useful addition bound.  We record it separately because
the answer is recursive rather than a single expression such as
$\lfloor n/2\rfloor+1$.

Here is the conclusion before the bookkeeping: the same circuits use
\[
  A_n\le
  \min\left\{2n,\ \frac54n+6\lceil\log_2 n\rceil^2+1\right\}
\]
additions or subtractions.  Thus the uniform constant is~$2$, while the
asymptotic leading constant is~$5/4$.  The first part of this section records
the exact, share-aware recurrences and then proves these two bounds
(the closing subsection treats multiplicative height);
a reader interested only in the headline estimate may skip directly to
\Cref{thm:construction-addition-count,thm:construction-addition-asymptotic}.

Let $\mathcal A$ count binary additions and subtractions of field values.
Negation, copying a wire, and multiplication by a fixed field constant are
free, as in the straight-line-program model of \Cref{sec:model}; every
data-dependent shift such as $U+\alpha$ costs one addition.  A let-bound
subexpression used more than once is charged only once.  Thus this is an
ordinary DAG count, not a count of plus signs after fully expanding a
formula.  If fixed integer multiples must themselves be synthesized from
additions, one may add the length of a chosen addition--subtraction chain at
each displayed multiple $(k-1)U$; this hardware-dependent variant is not the
cost model used here.

Write
\[
 q_1=1,\qquad q_s=5\,2^{s-2}-2\quad(s\ge2)
\]
for the addition count of $Q_{2^s-1}$, and
\[
 f_1=6,\qquad
 f_s=5(2^{s-1}+2^{s-2})-2\quad(s\ge2)
\]
for that of the full fill $A_{2^s}$.  The $s\ge2$ formulas are
\Cref{lem:fill-Q-count}; the two omitted direct bases are
$Q_1=x+\alpha_0$ and
\[
 A_2=(x+\beta_0)\bigl((H_2+\beta_1)S^{(1)}+\alpha_1\bigr)
       +(H_2+\beta_2)S^{(2)}+\alpha_0.
\]

There are two small sharing improvements in the shared $T$ bases.  Every
costed level-one call is supplied as
$\widetilde H_2=H_2+\rho$, so the circuit retains $\rho$ instead of
materializing $\widetilde H_2$ and later subtracting $H_2$ again.  It forms
$\widetilde H_4=H_4+\rho$ directly.  At the following shared odd base, put
\[
 F_1=H_4-(k-1)U_1.
\]
Since $U_2=U_1-\rho$ and $\widetilde H_4=H_4+\rho$, the second factor is
not formed independently:
\begin{equation}
 \widetilde H_4-(k-1)U_2=F_1+k\rho.
 \label{eq:shared-factor}
\end{equation}
Both rewrites preserve the construction identities and all decoder pivots;
they only remove redundant affine gates.  The proof-only quantity
$W_2=W_1+z$ is likewise not materialized, since the circuit directly forms
$\widetilde H_8=H_8+z$.

Let $\tau(k,l)$ be the additions used by this shared schedule for
$(T^{(1)}_{k,2^l},T^{(2)}_{k,2^l})$, with the power inputs already
available (and with the scalar shift retained at the two shared bases).
The algorithms give the exact recurrence
\begin{align}
 \tau(1,l)&=0, \label{eq:T-add-one}\\
 \tau(2m,1)&=\tau(m,2)+5, \label{eq:T-add-even-base}\\
 \tau(2m,l)&=\tau(m,l+1)+2q_{l-1}+8
     &&(l\ge2), \label{eq:T-add-even}\\
 \tau(2m+1,2)&=\tau(m,3)+15
     &&(m\ge1), \label{eq:T-add-odd-base}\\
 \tau(2m+1,l)&=\tau(m,l+1)+q_{l-1}+2q_{l-2}+q_l+16
     &&(m\ge1,\ l\ge3). \label{eq:T-add-odd}
\end{align}
For example, the $5$ in \eqref{eq:T-add-even-base} consists of four
additions for
$(H_2+(x+a))(H_2-(x+a))+e$ and one for $H_4+\rho$.
In \eqref{eq:T-add-odd-base}, the $Q_3$ call costs three, the two affine
inputs $U_1,V_1$ cost three, the common octic costs four, the shifted octic
costs one, and the two factors and two tails cost four by
\eqref{eq:shared-factor}, for a total of fifteen.
The constants $8$ and $16$ in the ordinary branches are obtained by the
same literal count of their two difference-of-squares products and output
assemblies.  Hence the recurrence charges every displayed gate exactly once.

For reference, let $g_d$ denote the addition count of the selected
degree-$d$ auxiliary gadget $\mathcal Q_d$.  The three gadget constructions
give
\begin{align}
 g_1&=1,&g_3&=3,&g_7&=8,\label{eq:g-add-bases}\\
 g_{4k+1}&=\tau(2k,1)+3 &&(k\ge1),\label{eq:g-add-4k1}\\
 g_{8k+3}&=\tau(2k,2)+9 &&(k\ge1),\label{eq:g-add-8k3}\\
 g_{8k+7}&=\tau(k,3)+19 &&(k\ge1).\label{eq:g-add-8k7}
\end{align}
Indeed, \eqref{eq:g-add-4k1} charges one addition for
$\widehat H_2=H_2+\gamma$ and two for
$(x+\beta)S^{(1)}+S^{(2)}$; the scalar difference is passed directly to
the shared base.  The known-powers construction in full generality costs
\[
 q_{l-1}+2+\tau(2k,l)+f_{l-1}
 \tag{known-powers-additions}
\]
for degree $2^{l+1}k+(2^l-1)$; setting $l=2$ gives
\eqref{eq:g-add-8k3}.  Finally, the barred gadget uses five additions for
$H_8$, one for $\widetilde H_8$, thirteen for $A_4$, and
$\tau(k,3)$ internally, proving \eqref{eq:g-add-8k7}.
The low slot in the $8k+3$ construction has one deliberate exception:
when $k=1$ the algorithm uses the scalar $\alpha_1$, not the polynomial
$Q_1=x+\alpha_1$.  Accordingly, put
\[
  g_d^{\mathrm{lo}}=g_d\quad(d\ge3),
  \qquad g_1^{\mathrm{lo}}=0.
\]

Let $a_n$ be the additions used to jointly produce the compatible pair in
odd degree $n\ne7$.  The finite ledgers and the three induction branches are
then
\begin{equation}
 \begin{gathered}
 a_3=3,\qquad a_{15}=23,\\
 a_{27}=43,\qquad a_{31}=43.
 \end{gathered}
 \label{eq:pair-add-bases}
\end{equation}
\begin{align}
 a_{4k+1}&=\tau(2k,1)+2,
 \label{eq:pair-add-4k1}\\
 a_{8k+3}&=a_{2k+1}+g_{4k+1}+g^{\mathrm{lo}}_{2k-1}+7,
 \label{eq:pair-add-8k3}\\
 a_{8k+7}&=a_{2k+1}+g_{2k+1}+g_{4k+3}+6.
 \label{eq:pair-add-8k7}
\end{align}
The first two recursive lines apply for $k\ge1$ (with $k\ne3$ in the
$8k+3$ line), and the last for $k\ge2$, $k\ne3$.
Here $a_{4k+1}$ includes the two additions forming $H_2$; the scalar
$\rho$ is passed to the shared $T$ base without first forming
$\widetilde H_2$.  The constants $7$ and $6$ are the literal costs of the
two outer difference-of-squares assemblies in the $8k+3$ and $8k+7$
algorithms.  In the latter, the two common forms
$S^{(1)}_3\mathbin{\pm}S^{(1)}_2$ are reused and the two scalar coordinates
are $s=a+b,d=b-a$.  This fixed linear reparameterization is invertible under
the admissibility hypothesis, so it does not change decodability.  For the
finite cases, the corresponding ledgers are
\[
\begin{array}{c|l|c}
 n&\text{addition costs, with shared producers charged once}&a_n\\ \hline
 3&H_2:2,\ T^{(2)}=H_2+\alpha_0:1&3\\
 15&H_2:2,\ H_4:4,\ Q_7:8,\ \text{outer assemblies}:9&23\\
 27&H_2:2,\ Q_{13}:23,\ Q_3:3,\ Q_7:8,\
       \text{outer assemblies}:7&43\\
 31&H_2,H_4:6,\ \bar Q_{15}:19,\ Q_7:8,\ Q_3:3,\
       \text{outer assemblies}:7&43.
\end{array}
\]

Finally let $A_n$ count additions in the complete degree-$n$ polynomial.
At top level in degree three we use the direct $Q_3$ circuit with
$H_2=x^2$; it costs three additions, one fewer than first producing the
degree-three pair and then combining it.  (The pair is still used inside
the recursion.)  The displayed septic circuit has eleven additions, but an
invertible affine change of its first three coordinates saves one.  Namely,
with fresh coordinates $\beta_i$, replace its last line by
\[
 P^{\rm add}_7=\beta_0+y+
   (\beta_2+x+z)(\beta_1+w).
\]
This is exactly the septic of \Cref{lem:septic-base} after the substitution
\[
 \alpha_0=\beta_0+\beta_1,\qquad
 \alpha_1=\beta_1,\qquad
 \alpha_2=\beta_2-1,
 \qquad \alpha_i=\beta_i\ (3\le i\le6).
\]
Thus its explicit decoder, followed by
$\beta_0=\alpha_0-\alpha_1$, $\beta_1=\alpha_1$,
$\beta_2=\alpha_2+1$, proves decodability.  The new last line costs two
additions instead of three, so the optimized septic costs ten.  In every
other odd pair case the final combination costs one addition, as does the
even lift, so
\begin{equation}
 \begin{aligned}
 A_1&=1,\qquad A_2=2,\qquad A_3=3,\qquad A_7=10,\\
 A_n&=a_n+1 &&(n\text{ odd},\ n\notin\{3,7\}),\\
 A_n&=A_{n-1}+1 &&(n\ge4\text{ even}).
 \end{aligned}
 \label{eq:full-add-recurrence}
\end{equation}
Equations
\eqref{eq:T-add-one}--\eqref{eq:full-add-recurrence} are an exact,
share-aware addition count for every construction in the theorem.

\begin{theorem}[Uniform addition bound]
\label{thm:construction-addition-count}
The optimized schedules above compute every complete degree-$n$
construction with at most $2n$ additions or subtractions.  More precisely,
$a_n\le2n-1$ for every costed odd compatible pair, and $A_n\le2n$ for
every $n\ge1$.
\end{theorem}
\begin{proof}
We include the short numerical induction so that the bound does not rest on
experimental enumeration.  Substituting the formulas for $q_s$ into
\eqref{eq:T-add-even}--\eqref{eq:T-add-odd} and inducting strongly on $k$
gives
\begin{align}
 \tau(k,2)&\le8(k-1)+2,\label{eq:T-add-bound-2}\\
 \tau(k,l)&\le2(k-1)2^l-\bigl(3\,2^{l-2}-4\bigr)
       &&(k\ge2,\ l\ge3),\label{eq:T-add-bound-high}\\
 \tau(2k,1)&\le8k-1.\label{eq:T-add-bound-1}
\end{align}
For clarity, the fresh nonrecursive costs used in that induction are
$5,10,5\,2^{l-2}+4$ in the three even cases
$l=1,l=2,l\ge3$, and $15,29,5\,2^{l-1}+8$ in the odd cases
$l=2,l=3,l\ge4$.  These are respectively no larger than the difference
between the right-hand side at $(k,l)$ and the right-hand side at the
smaller recursive argument.  This proves the three displayed inequalities;
at $l\ge3$ the worst boundary case is $k=2$, where
\eqref{eq:T-add-bound-high} is an equality.

Splitting $k$ by parity, and once more by its residue modulo four when
needed to select \eqref{eq:g-add-8k3} or \eqref{eq:g-add-8k7}, the same
three bounds give
\begin{align}
 g_{4k+1}+g^{\mathrm{lo}}_{2k-1}&\le12k-3 &&(k\ge1),
 \label{eq:g-pair-bound-3}\\
 g_{2k+1}+g_{4k+3}&\le12k+3 &&(k\ge2).
 \label{eq:g-pair-bound-7}
\end{align}
Here are the only estimates used in that residue check.  If $k\ge3$ is odd,
$g_{4k+1}\le8k-1$ and
$g^{\mathrm{lo}}_{2k-1}\le4k-2$; if $k$ is even,
$g_{4k+1}\le8k+2$ and
$g^{\mathrm{lo}}_{2k-1}\le4k-5$.  The case $k=1$ is checked directly.
For the second line, writing $k=2r$ gives the exact upper ledger
\[
 g_{2k+1}+g_{4k+3}
 \le \tau(r,2)+\tau(r,3)+27\le24r+3=12k+3;
\]
when $k=2r+1$, the estimates
$g_{4r+3}\le8r+3$ and
$g_{8r+7}\le16r+1$ for $r\ge2$ give a stronger bound, and the case $r=1$
is direct.  Thus \eqref{eq:g-pair-bound-3}--\eqref{eq:g-pair-bound-7}
follow without a finite search.

We now use strong induction on the odd target degree.  The four values in
\eqref{eq:pair-add-bases} satisfy $a_n\le2n-1$.  For $n=4k+1$,
\eqref{eq:T-add-bound-1} and \eqref{eq:pair-add-4k1} give
$a_n\le8k+1=2n-1$.  In the $8k+3$ branch, the induction hypothesis and
\eqref{eq:g-pair-bound-3} give
\[
 a_{8k+3}\le(4k+1)+(12k-3)+7=16k+5=2(8k+3)-1.
\]
In the $8k+7$ branch, \eqref{eq:g-pair-bound-7} gives
\[
 a_{8k+7}\le(4k+1)+(12k+3)+6
     =16k+10<2(8k+7)-1.
\]
This proves the pair bound.  Adding the final combination gives
$A_n\le2n$ for odd $n\notin\{3,7\}$; the direct cubic and septic have
$A_3=3$ and $A_7=10$.  The other two small bases are immediate, and an
even lift adds only one gate, so
$A_n=A_{n-1}+1\le2(n-1)+1<2n$.  The theorem follows.
\end{proof}

The uniform factor-two statement is convenient for small degrees, but it
does not describe the asymptotic behaviour of the same schedules.  Their
leading addition constant is in fact $5/4$, not~$2$.

\begin{theorem}[Sharper asymptotic addition bound]
\label{thm:construction-addition-asymptotic}
Put
\[
  L(n)=\lceil\log_2 n\rceil,
  \qquad L(1)=0.
\]
The optimized schedules above satisfy
\[
  A_n\le \frac54n+6L(n)^2+1.
\]
Consequently
\[
  A_n\le
  \min\left\{2n,\ \frac54n+6\lceil\log_2n\rceil^2+1\right\},
\]
so they use at most $\frac54n+O((\log n)^2)$ additions.  In particular,
$A_n\le\frac32n$ for every $n\ge4096$.
\end{theorem}
\begin{proof}
We first sharpen the coarse estimates on the $T$ recursion.  The following
slightly stronger high-level estimate is convenient:
\begin{equation}
 \tau(k,l)\le
 5\,2^{l-2}(k-1)+8L(k)-4
 \qquad(k\ge2,\ l\ge4).
 \label{eq:T-add-sharp-high}
\end{equation}
This follows by strong induction on $k$.  Put $p=2^{l-2}$.  At an even
step the fresh cost is $5p+4$, while at an odd step it is $10p+8$.
The case $k=2$ is an equality.  If $k=2m$ with $m\ge2$, apply the
induction hypothesis to $(m,l+1)$ and use
$L(2m)=L(m)+1$; the resulting bound is four smaller than the right-hand
side of \eqref{eq:T-add-sharp-high}.  If $k=2m+1$, the case $m=1$ is
direct, and for $m\ge2$ the induction hypothesis together with
$L(2m+1)\ge L(m)+1$ gives the claimed bound directly.

The exceptional level $l=3$ has fresh costs $14$ and $29$.  Applying
\eqref{eq:T-add-sharp-high} at level four (and checking $k=1,2,3$
directly) gives
\begin{equation}
 \tau(k,3)\le10(k-1)+8L(k)
 \qquad(k\ge1).
 \label{eq:T-add-sharp-three}
\end{equation}
For $l=2$, use
\[
 \tau(2m,2)=\tau(m,3)+10,
 \qquad
 \tau(2m+1,2)=\tau(m,3)+15.
\]
Together with $L(2m)=L(m)+1$ and
$L(2m+1)\ge L(m)+1$, these give the same bound
\begin{equation}
 \tau(k,2)\le5(k-1)+8L(k).
 \label{eq:T-add-sharp-two}
\end{equation}

Substitution in \eqref{eq:g-add-bases}--\eqref{eq:g-add-8k7} now gives one
uniform estimate for every selected odd gadget:
\[
  g_d\le\frac54d+8L(d).
  \tag{gadget-sharp-bound}
\]
For example,
$g_{4k+1}=\tau(k,2)+8$,
$g_{8k+3}=\tau(2k,2)+9$, and
$g_{8k+7}=\tau(k,3)+19$; the three direct bases are immediate.

We next prove
\[
  a_n\le\frac54n+6L(n)^2
  \tag{pair-sharp-bound}
\]
by strong induction over the costed odd pairs.  The four finite bases satisfy
it directly.  If $n=4k+1$, then
\[
 a_n=\tau(k,2)+7
   \le5(k-1)+8L(k)+7
   \le\frac54n+6L(n)^2.
\]
In either of the remaining branches let $m=2k+1$ be the smaller pair and
let $d_1,d_2$ be the two auxiliary-gadget degrees (with the degree-one low
slot omitted when $k=1$).  In both cases
\[
 d_1+d_2\le n-m-2,
 \qquad L(m)\le L(n)-2.
\]
The outer assembly costs at most seven additions.  Writing $L=L(n)$, the
induction hypothesis and (gadget-sharp-bound) therefore give
\[
 \begin{aligned}
  a_n
  &\le \frac54m+6(L-2)^2
       +\frac54(d_1+d_2)+16L+7\\
  &\le \frac54n+6(L-2)^2+16L+\frac92
   \le \frac54n+6L^2.
 \end{aligned}
\]
The last inequality is $8L\ge57/2$; every recursive target here has
$L\ge4$.  This proves (pair-sharp-bound).

For an odd complete construction, the last combination adds one, except
for the already smaller direct cubic and septic bases.  An even lift also
adds one and satisfies $L(n-1)\le L(n)$.  Hence
$A_n\le\frac54n+6L(n)^2+1$ for every $n$.  Combining this with
\Cref{thm:construction-addition-count} proves the minimum displayed in the
statement.  Finally,
$6L(n)^2+1\le n/4$ for $n\ge4096$ (check the first two dyadic boundary
points; thereafter the right-hand side doubles while the left grows
quadratically in $L$), which gives the last assertion.
\end{proof}

\begin{remark}[What would be needed for an $n+o(n)$ addition bound]
The remaining quarter is structural in the present schedule, not an artefact
of the final estimate: already the Mersenne block has the exact count
\[
 \mathcal A(Q_{2^s-1})=5\,2^{s-2}-2
   =\frac54(2^s-1)-\frac34.
\]
The affine sharing above removes repeated shifts and outer assembly gates,
but it does not change this leading term.  Reaching $n+o(n)$ additions while
retaining the exact $\lfloor n/2\rfloor+1$ multiplication count would
therefore require a genuinely more addition-efficient fill (or a different
known-powers gadget), rather than a sharper analysis of the current DAG.
Rabin and Winograd obtain $n+o(n)$ additions when an additional
$O(\log n)$ multiplications are allowed~\cite{rabin1972number}; eliminating
both overheads simultaneously remains open here.
\end{remark}

\subsection*{Multiplicative height}

Gate counts do not describe the parallel (or pipelined) cost of the same
schedules, so we also record their \emph{height}: the maximum number of
nonscalar multiplications on a directed path from an input to the output of
the chain, in the same cost model as above (additions, subtractions,
negations and fixed scalar multiples are free).  Horner's rule has height
$n-1$, and the hand-optimized circuits of the appendix have height close to
$\log_2n$.  The shipped schedules match this: the peeled known-powers
gadget of \Cref{sec:peeled-Q} gives height $2\lceil\log_2n\rceil+4$ for odd
$n$ and $2\lceil\log_2n\rceil+5$ for even $n$ at the identical gate ledger
(\Cref{thm:construction-height-peeled}).  The
original sequential-fill scheduling, analyzed below as the baseline, has
height $\Theta((\log n)^2)$, with a small leading constant.

\begin{theorem}[Height of the sequential baseline schedules]
\label{thm:construction-height}
Let $L=L(n)=\lceil\log_2n\rceil$.  The sequential-fill schedules above compute
every complete degree-$n$ construction with height at most
\[
  \tfrac14L^2+L+2 \quad (n\text{ odd}),\qquad
  \tfrac14L^2+L+3 \quad (n\text{ even}).
\]
\end{theorem}
\begin{proof}
Write $h(U)$ for the height of a wire $U$, so $h(x)=0$ and every displayed
product raises the maximum operand height by exactly one.  The proof is a
ledger induction parallel to \Cref{thm:construction-addition-count}, with
four invariants.

\emph{Fill.}  The first component of the fill $A_{2^\ell}$
(\Cref{alg:constr-fill}) multiplies $S^{(1)}$ by one factor at each level
$\ell,\ell-1,\dots,1$ and finishes with the outer $(x+\beta_0)$ gate, a
sequential chain of $\ell+1$ products; the factor consumed at level $j$ is
$H_{2^j}+Q_{2^{j-1}-1}$, and the second component is the same chain with
scalar-shifted factors.  Given the tower and known-powers invariants below
for all indices $\le\ell$, this yields
\begin{equation}
  h\bigl(A_{2^\ell}(S^{(1)},S^{(2)})\bigr)
  \le\max\Bigl(h(S^{(1)}),\,h(S^{(2)}),\,
    \bigl\lfloor\tfrac{\ell^2}{4}\bigr\rfloor+1\Bigr)+\ell+1.
  \label{eq:fill-height}
\end{equation}

\emph{Tower and known powers.}  By simultaneous induction,
\begin{equation}
  h(H_{2^i})\le\bigl\lfloor\tfrac{i^2}{4}\bigr\rfloor+1
  \quad(i\ge1),
  \qquad
  h(Q_{2^t-1})\le\bigl\lfloor\tfrac{(t+1)^2}{4}\bigr\rfloor
  \quad(t\ge1).
  \label{eq:tower-height}
\end{equation}
Indeed, every displayed power step forms $H_{2^{\ell+1}}$ with one new
product whose operands involve $H_{2^\ell}$ and blocks of index at most
$\ell-1$, so
$h(H_{2^{\ell+1}})
  \le1+\max\bigl(\lfloor\ell^2/4\rfloor+1,\lfloor\ell^2/4\rfloor\bigr)
  =\lfloor\ell^2/4\rfloor+2
  \le\lfloor(\ell+1)^2/4\rfloor+1$
for $\ell\ge2$, and the two shared bases are checked directly.  For the
known-powers gadget at $t\ge4$,
$Q_{2^t-1}=A_{2^{t-2}}\bigl(H_{2^{t-1}}+Q_{2^{t-2}-1},\,
H_{2^{t-1}}+\alpha\bigr)$, so \eqref{eq:fill-height} gives
$h(Q_{2^t-1})\le\bigl(\lfloor(t-1)^2/4\rfloor+1\bigr)+(t-1)
  =\lfloor(t+1)^2/4\rfloor$,
and $t\le3$ is direct.  The tilde towers and the good-polynomial gadgets
$\bar Q$ use the same fills and one square-difference per level with
operands of index at most $\ell$, so they satisfy the same two bounds.

\emph{The $T$ recursion.}  For $T_{k,2^\ell}$ put
$\lambda=\ell+\lceil\log_2k\rceil$ and let $s(k)\le\lceil\log_2k\rceil$ be
the number of odd values exceeding $1$ in the halving chain of $k$.  Then
\begin{equation}
  h\bigl(T^{(1)}_{k,2^\ell}\bigr),\,
  h\bigl(T^{(2)}_{k,2^\ell}\bigr)
  \le\bigl\lfloor\tfrac{\lambda^2}{4}\bigr\rfloor+1+s(k).
  \label{eq:T-height}
\end{equation}
The base $k=1$ is the tower bound; an even step recurses at $(k/2,\ell+1)$
with the same $\lambda$ after one tower step; an odd step multiplies the
recursive pair by factors of height at most
$\lfloor(\ell+1)^2/4\rfloor\le\lfloor\lambda^2/4\rfloor+1+s\bigl(\tfrac{k-1}2\bigr)$
and adds a $Q$ block, contributing exactly the $+1$ that separates $s(k)$
from $s\bigl(\tfrac{k-1}2\bigr)$.

\emph{Assembly.}  Let $B(L)=\lfloor L^2/4\rfloor+L+1$.  By strong induction
on the odd degree, every costed pair of degree $m$ has height at most
$B\bigl(L(m)\bigr)$: the $4k+1$ family is \eqref{eq:T-height} with
$\lambda=1+\lceil\log_2 2k\rceil\le L$ for $T_{2k,2}$ and $s\le\lambda-1$, giving height at most $\lfloor L^2/4\rfloor+L<B(L)$; the $8k+3$ and $8k+7$ steps add one
square-difference on top of $Q$ or $\bar Q$ blocks at level $\le L-1$
(height $\le\lfloor L^2/4\rfloor+1$ by \eqref{eq:tower-height}) and of a
recursive pair of degree $\le(n-1)/2$ (height $\le B(L-1)$); and the
special cases $15$, $27$, $31$ are checked directly.  The final combination
$P_n=x\cdot T^{(1)}+T^{(2)}$ adds one product, giving
$B(L)+1=\lfloor L^2/4\rfloor+L+2\le\frac14L^2+L+2$.  For even $n$ the
lift $c_0+xQ_{n-1}$ of \Cref{thm:construction-count} adds one further
product on top of the odd-degree schedule, and $L(n-1)\le L(n)$.
\end{proof}

\begin{remark}[The order $(\log n)^2$ is genuine]
The quadratic term comes from the sequential fill chains: level $\ell$ of
the tower waits for a fill of length $\Theta(\ell)$ whose factors wait for
the tower below it.  The bound reflects the schedules themselves, not slack
in the analysis: the exact heights (reported by
\texttt{tools/polychain.py}, whose \texttt{selftest} subcommand also
checks the bound of \Cref{thm:construction-height} for every
$3\le n\le N$ with $N=120$ by default; the heights quoted here were
computed for $3\le n\le6001$) grow as
$\Theta((\log n)^2)$, e.g.\ height $5$ at $n=17$, $10$ at $n=63$, $12$ at
$n=127$, $16$ at $n=253$, $25$ at $n=1021$ and $36$ at $n=4093$; the
observed heights over $3\le n\le6001$ never exceed
$\lfloor L^2/4\rfloor+3$ (attained only at the even lifts $n=8,16,32$; for
$7\le L\le12$ the per-$L$ maximum is $\lfloor L^2/4\rfloor+1$, attained at
$n=2^L-2$), and the best case (near $n=2^{L-1}+3$) is
$\approx L^2/6$.  As with the addition constant, this is a property of the
fill.

Depth-balanced decodable known-powers gadgets exist at no gate cost at
all: the peeled recursion \eqref{eq:peeled-Q} matches the ledger of
\Cref{lem:fill-Q-count} exactly, with height $k$ in place of the quadratic
growth, and substituting it at every known-powers site --- leaving the
keyed tower spine, seam fills and good-polynomial gadgets byte-for-byte
unchanged, so the multiplication count remains exactly
$\lfloor n/2\rfloor+1$ and no Rabin--Winograd-style key-free squarings
appear --- removes the quadratic term entirely; see \Cref{sec:peeled-Q}, \Cref{fig:peeled-Q},
and \Cref{thm:construction-height-peeled}.
\end{remark}

\begin{corollary}[All odd degrees have a decodable construction]
\label{cor:all-odd-decodable}
Over every $n$-admissible field, each odd $n\ge3$ has a decodable monic
degree-$n$ construction.  For $n=7$ use
\Cref{lem:septic-base}; for every other odd degree use
\Cref{thm:odd-realizable-pairs}.
\end{corollary}
